% Panel captions for musculo-skeletal-syntax.tex
%
% Every numeric value quoted below is read from the JSON records written by
% the reference implementation (musculo-skeletal/python/results/), not from
% the prose. Panels 1-8 report the executed programs; panels 9-10 report the
% coherence, identity, opacity and floor validation suite (seed 20260819,
% 2432 checks, all passing).

\begin{figure*}[!htbp]
    \centering
    \includegraphics[width=\textwidth]{figures/panel1-closure-dynamics.png}
    \caption{\textbf{Closure decides the character of the dynamics, not
    merely its magnitude.}
    The closure index of every arm below is decided statically from the
    declared paths in linear time (\Cref{thm:closure-decidable}), before any
    integration is performed, so the separation between the two families is
    fixed by the program text rather than discovered in the trace.
    (\textbf{A})~State of a corticospinal circulation over $12$\,s, closed
    (left axis) and after severance of the cortical output element (right
    axis). The two require separate scales because they differ by four orders
    of magnitude: the closed circulation executes bounded oscillation about a
    non-zero centre, as \Cref{cor:perpetual} requires of an ungrounded graph
    under sustained supply, while the severed one ramps monotonically until
    it saturates the integrator. No value of a gain interpolates between
    these.
    (\textbf{B})~Phase portrait of the closed circulation, state against its
    time derivative over $9$\,s. The trajectory fills an annular region and
    never approaches the marked centroid, which is the geometric content of
    \Cref{thm:no-equilibrium}: the system has no fixed point to relax onto,
    so what electrophysiology calls a resting potential is the centre of a
    cycle rather than a state the circuit occupies.
    (\textbf{C})~Excursion envelope on a logarithmic axis, closed and open,
    with the median divergence time across all open arms in the corpus marked
    at $3.72$\,s. The closed envelope is stationary near $10^{-2}$; the open
    envelope crosses it within the first second and saturates four decades
    higher. Divergence is reported as a time, not raised as an exception,
    which is the totality obligation a conforming backend must discharge.
    (\textbf{D})~Envelope of the closed circulation over loop gain and
    elapsed time, with level contours projected onto the base plane. Decay is
    governed by the product of gain and time, so a weakly conducting return
    sustains a wider limit cycle without ever failing to close: the surface
    approaches zero envelope only asymptotically and never leaves the closed
    regime.}
    \label{fig:panel1}
\end{figure*}

\begin{figure*}[!htbp]
    \centering
    \includegraphics[width=\textwidth]{figures/panel2-attenuation-severance.png}
    \caption{\textbf{Attenuation and severance are different in kind, and the
    language refuses to let a scaling factor imitate a cut.}
    The sweep below is exhaustive over twelve decades of $\lambda$ rather than
    sampled, so the invariance is a property of the analysis and not an
    artefact of the grid: \Cref{prop:attenuation} is checked, not illustrated.
    (\textbf{A})~Closure index against attenuation factor $\lambda$ down to
    $10^{-12}$. The index remains \texttt{closed} at every $\lambda>0$, while
    the single removal \texttt{without element(k)} opens the circuit outright.
    The two operators are therefore not endpoints of one continuum, and a
    program cannot reach an aperture by shrinking a gain.
    (\textbf{B})~Force output through the six declared phases of a peripheral
    nerve block: $420$, $420$, $420$, $0$, $126$, $420$\,N. Analgesia and
    proprioceptive loss leave force intact although the third phase is already
    open; motor block abolishes it; recovery restores $30\%$ before resolution
    returns the full value. Force reads the outbound phase alone, so an open
    return does not inflate it.
    (\textbf{C})~Divergence time against $\lambda$, showing the discontinuity
    at $\lambda=0$. Every $\lambda>0$ yields no divergence; the removal yields
    $3.72$\,s. The gap at the origin is the quantitative form of the design
    commitment.
    (\textbf{D})~Loop-gain surface over $\lambda$ and elapsed time. The
    surface descends smoothly toward the $\lambda\to0$ edge without ever
    reaching the open plane, which is why a limit argument on the attenuation
    parameter does not recover severance.}
    \label{fig:panel2}
\end{figure*}

\begin{figure*}[!htbp]
    \centering
    \includegraphics[width=\textwidth]{figures/panel3-lesion-level.png}
    \caption{\textbf{Upper- and lower-motor-neuron syndromes separate without
    either being encoded: the taxonomy falls out of which phase the lesion
    opens.}
    Neither program declares spasticity, flaccidity, nor a lesion level. Each
    declares only a circulation and an element to remove, so the clinical
    contrast below is derived rather than restated.
    (\textbf{A})~Tonic discharge rate by arm. The supraspinal circulation
    after a cortical lesion sustains no rhythm at all, while the segmental
    circulation in the same program sustains $15.72$\,Hz --- the descending
    path is severed, the local loop is not, and rhythmicity survives exactly
    where the loop does. The anterior-horn arm supports neither.
    (\textbf{B})~Loop latency across the arms: $52.8$, $31.8$, $60.8$,
    $38.8$ and $49.8$\,ms. Latency shortens when a lesion removes a long
    descending limb and lengthens when it removes a fast one, so latency alone
    is not a severity ordering.
    (\textbf{C})~Below- and above-lesion behaviour of a complete T6
    transection: reflex activity below the lesion persists at $12.89$\,Hz
    while voluntary function above it persists at $10.04$\,Hz, the two now
    uncoupled. That both survive is the point --- a transection is a partition
    of circulations, not a global loss of function.
    (\textbf{D})~State trajectories of the intact, upper-motor-neuron and
    lower-motor-neuron arms in a common space of activation, rate and time.
    The three occupy disjoint regions, and the separation is generated by the
    closure structure alone, with no syndrome-specific parameter anywhere in
    the source.}
    \label{fig:panel3}
\end{figure*}

\begin{figure*}[!htbp]
    \centering
    \includegraphics[width=\textwidth]{figures/panel4-spectral-separation.png}
    \caption{\textbf{Tremor type is read from which stratum carries the power,
    and the type-separation statistic reports honestly when it cannot be
    read.}
    Band powers are integrated over the stratum bands fixed in
    \Cref{def:strata}, which are declared in the specification and not tuned
    per arm, so no arm receives a band chosen to flatter it.
    (\textbf{A})~Reflex-band power across intact, parkinsonian, essential and
    cerebellar arms: $0.241$, $0.278$, $0.229$, $0.281$. The four are close,
    and reflex power alone does not discriminate.
    (\textbf{B})~Supraspinal-band power for the same arms: $0.0123$, $0.0188$,
    $0.0226$, $0.0198$. The parkinsonian and cerebellar arms carry roughly
    $1.5$--$1.6\times$ the intact supraspinal power and the essential arm
    nearly $1.8\times$; the ratio between strata, not either band alone,
    carries the taxonomy.
    (\textbf{C})~Cross-stratum coupling index by arm: $0.052$, $0.048$,
    $0.063$, $0.047$. The essential arm is the outlier, consistent with a
    lesion that widens coupling rather than shifting power between strata.
    (\textbf{D})~Type separation $\eta$ over band power and coupling. Three
    arms attain $\eta \geq 0.899$, but the essential arm attains only $0.166$.
    Under \Cref{prop:degeneracy} this is the estimator declaring that its
    types are not separated for that arm, so the estimate is reported together
    with the warning rather than as a number of equal standing --- the
    \emph{unconditional} claim that a typed estimate is available for every
    arm is false, and \Cref{rem:report-eta} requires that this be stated
    rather than absorbed.}
    \label{fig:panel4}
\end{figure*}

\begin{figure*}[!htbp]
    \centering
    \includegraphics[width=\textwidth]{figures/panel5-repair-closure.png}
    \caption{\textbf{Surgical repair is re-closure, and its cost is a latency
    that the closure index by itself cannot see.}
    All four arms below are read from one program in which only the declared
    paths differ; the reroute operator is the sole construct in the language
    permitted to raise a closure index, so a repair is syntactically visible
    as such.
    (\textbf{A})~Closure index by arm. The amputated arm is open; targeted
    muscle reinnervation and the mirror construction both restore
    \texttt{closed}. On the closure axis alone the two repairs are
    indistinguishable from the intact circuit, which is precisely the
    limitation that motivates the remaining three views.
    (\textbf{B})~Loop latency: intact $49.8$\,ms, TMR $65.8$\,ms ($+16.0$),
    mirror $35.8$\,ms ($-14.0$). Re-closure through a rerouted path is not
    latency-neutral, and the two repairs err in opposite directions --- TMR
    routes through a longer path, the mirror through a shorter one.
    (\textbf{C})~Latency deviation from intact against restored closure,
    showing that equal closure admits a $30$\,ms spread. Closure is a
    predicate; the cost of achieving it is not recoverable from the predicate.
    (\textbf{D})~State trajectories of the four arms in activation, rate and
    time. The amputated arm departs the shared region and does not return; the
    two repaired arms re-enter it along visibly different approaches, so the
    difference the closure index cannot express is present in the dynamics.}
    \label{fig:panel5}
\end{figure*}

\begin{figure*}[!htbp]
    \centering
    \includegraphics[width=\textwidth]{figures/panel6-coupled-stiffness.png}
    \caption{\textbf{Joint stiffness is a property of a coupled antagonist
    pair, not of either muscle, and losing reciprocal inhibition raises it.}
    Coupling is counted only over compartments lying on both declared loop
    paths, so the shared-vertex count is a structural fact about the program
    rather than a name collision between two independently declared circuits.
    (\textbf{A})~Co-contraction ratio for the elbow pair: $0.6258$ intact
    against $0.8682$ with reciprocal inhibition removed, a rise of $38.7\%$.
    The direction matters --- removing the inhibition that separates the
    antagonists must increase simultaneous activation, and an earlier
    formulation that reported the opposite was diagnosed by exactly this
    check.
    (\textbf{B})~Derived joint stiffness: $138.58$ against
    $148.02$\,N\,m$^{-1}$, a $6.8\%$ increase. Stiffness is computed from the
    pair and is undefined for a single circulation, so it is an observable the
    antagonist-pair construct makes available and the single-loop language
    cannot express.
    (\textbf{C})~Oscillation amplitude, $0.01368$ against $0.01385$. The
    stiffer pair oscillates slightly more, not less: the added stiffness comes
    from co-activation rather than from damping, a distinction a lumped
    stiffness parameter would erase.
    (\textbf{D})~Stiffness surface over agonist and antagonist drive, with the
    intact and disinhibited operating points marked. The surface rises along
    the diagonal and is flat along either axis alone, which is the geometric
    statement that stiffness responds to the \emph{pair} and is insensitive to
    either member driven by itself.}
    \label{fig:panel6}
\end{figure*}

\begin{figure*}[!htbp]
    \centering
    \includegraphics[width=\textwidth]{figures/panel7-degenerate-estimator.png}
    \caption{\textbf{The instance-specific estimator agrees with its own
    measurement on data generated by no process at all.}
    Each generating process below contributes $4000$ independently drawn
    cascades, including one constructed specifically to violate the
    composition law, so an agreement that survives the adversarial corpus
    cannot be a property of the data.
    (\textbf{A})~Composed prediction against directly measured net power for
    the three processes: one obeying the composition law, one violating it,
    one a random walk. Every point of every process lies on the identity line,
    with Pearson $r = 1.000000$ throughout. Were the two quantities
    independent, the adversarial and random-walk clouds would scatter.
    (\textbf{B})~Maximum absolute discrepancy per process on a logarithmic
    axis, with double-precision epsilon marked: $2.2\times10^{-16}$,
    $1.8\times10^{-15}$, $1.4\times10^{-15}$. All three sit within one order
    of magnitude of machine epsilon. A statistic reporting perfect agreement
    on a random walk is not measuring the random walk (\Cref{cor:no-null}).
    (\textbf{C})~Residual over cascade length and process index. The surface
    is flat in the process direction --- agreement is insensitive to what
    generated the data, the signature of an algebraic identity rather than a
    confirmation (\Cref{thm:telescoping}) --- and rises only with length, as
    accumulated floating-point error.
    (\textbf{D})~Root-mean-square error of four candidate composition laws
    under the typed estimator on multiplicatively generated mixed-type
    cascades. The multiplicative law attains $0.0338$ against $0.191$, $0.333$
    and $0.452$; the typed test therefore selects among competitors, which is
    what a test possessing a null hypothesis is supposed to do
    (\Cref{thm:nondegenerate}). The corpora also bracket the condition:
    $\eta = 0.934$ on separated types against $\eta = 0.013$ on compressed
    ones, at nearly equal residual.}
    \label{fig:panel7}
\end{figure*}

\begin{figure*}[!htbp]
    \centering
    \includegraphics[width=\textwidth]{figures/panel8-capacitance-force.png}
    \caption{\textbf{Prosthetic compliance enters as a capacitance, and the
    force it yields is fixed by a ratio the program never states.}
    Charge is computed compartment-wise as $Q=\sqrt{2CP}$
    (\Cref{eq:charge}); mixing capacitances across compartments is rejected by
    \Cref{def:comp-rule}, so the two limbs below are compared, never summed.
    (\textbf{A})~Peak force at the biological and prosthetic hip: $474.5$ and
    $335.6$\,N. The ratio is $0.707107$, which is $\sqrt{0.9/1.8}$ to six
    figures --- exactly the square root of the declared capacitance ratio, and
    a consequence of \Cref{eq:charge} rather than a fitted coefficient.
    (\textbf{B})~Loop latency at both hips, $72.8$\,ms each. Compliance
    changes what force the limb can develop without changing when it develops
    it, so a latency-based gait metric is blind to the asymmetry panel A
    reports.
    (\textbf{C})~Oscillation amplitude, $0.01404$ and $0.01416$, and derived
    floor value $4.0$ at both limbs. The floor is derived from connectivity
    and positive edge weights (\Cref{prop:floor}) and is unchanged by the
    substitution, confirming that a prosthesis alters capacitance without
    altering the separation structure.
    (\textbf{D})~Force surface over compartment capacitance and supplied
    power, with both hips marked. The surface is a square root in each
    argument, so the two operating points lie on one sheet: no
    prosthesis-specific term is introduced, and the asymmetry is entirely a
    displacement along the capacitance axis.}
    \label{fig:panel8}
\end{figure*}

\begin{figure*}[!htbp]
    \centering
    \includegraphics[width=\textwidth]{figures/panel9-coherence-closure.png}
    \caption{\textbf{Closure and coherence are independent axes: a circuit can
    match its paths perfectly and still disagree with itself about them.}
    Coherence is computed from declared delays alone by signing each transport
    with its traversal direction, so a route taken backwards discharges the
    delay it accrues forwards and a cycle sums a \emph{disagreement} rather
    than a total. The margin is therefore backend-independent in the sense of
    \Cref{prop:backend-independent}, verified across two seeds on $60$
    circuits with agreement below $10^{-15}$.
    (\textbf{A})~Coherence margin against closure index over $200$ randomly
    generated circulations. All $200$ are closed, yet $76$ carry a strictly
    negative margin. Closure asks whether the declared paths \emph{match};
    coherence asks whether they \emph{agree}, and the first does not imply the
    second.
    (\textbf{B})~Distribution of route disagreement among the incoherent
    circuits, ranging from $4.9$ to $29.4$\,ms with median $16.8$\,ms. The
    disagreement is of the same order as the loop latencies themselves, so it
    is not a rounding-scale effect.
    (\textbf{C})~Margin before and after two expressible operators applied to
    $120$ incoherent circuits. Removal of the offending element repairs the
    disagreement in $120/120$ cases; scaling that element to $\lambda=0.05$
    leaves it incoherent in $120/120$. Both operators are expressible, so
    expressibility alone does not discriminate --- this is why admissibility
    is two-factor. Note also that scaling changes the margin by less than
    $10^{-12}$ and therefore \emph{passes} a coherence-preserving test while
    leaving the circuit broken: a preservation test cannot substitute for
    reporting the post-state margin.
    (\textbf{D})~Disagreement surface over route delay mismatch and the number
    of parallel routes. The surface is identically zero along the
    single-route edge and rises with mismatch only once a second route exists,
    which is the structural statement that incoherence requires an alternative
    and cannot arise on a circuit with one path.}
    \label{fig:panel9}
\end{figure*}

\begin{figure*}[!htbp]
    \centering
    \includegraphics[width=\textwidth]{figures/panel10-identity-opacity-floor.png}
    \caption{\textbf{Identity and character are two orientations of one cut,
    the interior of a path is not determined by its endpoints, and the floor
    bound is positive without being tight.}
    Minimum cuts are exact, obtained by enumeration over all partitions up to
    the declared vertex cap rather than by a heuristic solver, so no agreement
    below can be a solver artefact; where the cap binds, the analysis reports
    that it has fallen back rather than returning a silently approximate cut.
    (\textbf{A})~Minimum cut cost against cheapest singleton cut for two
    declared circuits. The serial chain gives $4.000$ against $4.000$ --- the
    cheapest cut \emph{is} a singleton, and one compartment carries the
    character. The bridged pair gives $2.000$ against $5.000$, a ratio of
    $2.50$: the minimising partition is a block cut into two groups of three,
    so no single compartment carries the character, and reading identity off
    any one compartment is impossible in principle rather than merely
    difficult.
    (\textbf{B})~Character cost of $120$ circuits against the same circuits
    under a random weight-preserving relabelling of every compartment. All
    $120$ lie exactly on the identity line, maximum deviation $0$ --- not
    small, but zero. Character is what survives relabelling; identity names
    particular compartments and does not survive it.
    (\textbf{C})~Distinct realisable interiors, shortest route latency and
    latency spread over $120$ circuits. $79$ of $120$ are opaque: endpoints
    and closure index admit more than one interior, with spreads reaching
    $19.1$\,ms and averaging $5.1$\,ms. An observer who records only where a
    path begins and ends, and whether it closes, has not thereby determined
    what it traverses.
    (\textbf{D})~Distribution of floor tightness, the ratio of separation cost
    to minimum edge weight, over $400$ circuits. The bound of
    \Cref{prop:floor} holds in every case, minimum $2.000$ and mean $2.502$,
    but the minimum is twice the bound rather than equal to it. The floor is
    therefore \emph{positive} by connectivity and positive edge weights, and
    \emph{not} tight: \Cref{prop:floor} is a guarantee, not an estimate, and
    \Cref{rem:floor-estimators} should be read accordingly.}
    \label{fig:panel10}
\end{figure*}
